%%%%%%%%%%%%
%%%%%%%%%%%%%%%%%%%%%%%%%%%%%%%%%%%%%%%%
%
% $ Autor: Gruppe 07 $
% $ Projekt: Magic Faucet Fountain $
% $ Pfad: LaTeXTemplate/Nano33BLESense/Demontageanleitung/Chapters/de/Lagerung.tex $
% $ Kurzbeschreibung: Dieses Dokument zeigt, wie der Magic Faucet Fountain gelagert wird. $
%
% !TeX spellcheck = en_GB/de_DE
% !TeX encoding = utf8
% !TeX root = manual 
% !TeX TXS-program:bibliography = txs:///biber
%
%%%%%%%%%%%%%%%%%%%%%%%%%%%%%%%%%%%%%%%%

\chapter{Lagerung}
Nach der Demontage sollten alle Bauteile sorgfältig getrocknet, gereinigt und in geeigneten Behältern oder Aufbewahrungslösungen verstaut werden. Achten Sie darauf, dass keine Feuchtigkeit zurückbleibt, um Schimmelbildung zu vermeiden.

Empfohlen wird die Verwendung kleiner, beschrifteter Behälter oder Beutel für empfindliche oder kleinteilige Komponenten wie Schrauben. Größere Elemente wie der Blumentopf, der Deckel oder die Steine können separat gelagert werden – möglichst staubgeschützt und trocken.

Falls Sie den Magic Faucet Fountain zu einem späteren Zeitpunkt wieder zusammenbauen möchten, empfiehlt es sich, die Komponenten sortiert und in der ursprünglichen Anordnung zu lagern. So vermeiden Sie Verwechslungen und vereinfachen die spätere Montage.